\subsection*{连续种群模型的平衡点}
对定义~\ref{def:continuous-logistic} 中的方程，现在可以使用 Lyapunov 稳定性与渐近稳定性的定义。
一维自治方程 $x'=f(x)$ 在相邻平衡点之间的符号决定运动方向；其右端在这里光滑，故解的唯一性由 Picard 定理保证。
\begin{proposition}[捕捞阈值]\label{prop:harvest-threshold}
定额捕捞模型在 $0<h<rK/4$ 时有两个正平衡点
\[
 x_\pm=\frac K2\left(1\pm\sqrt{1-\frac{4h}{rK}}\right).
\]
$x_-$ 不稳定，$x_+$ 渐近稳定。在临界值 $h=rK/4$ 时两者合并为
$K/2$；$h>rK/4$ 时没有平衡点，种群数量从任意正初值单调下降并在有限时间达到零。
比例捕捞模型在 $q<r$ 时有正的渐近稳定平衡点 $K(1-q/r)$；
$q\geq r$ 时每个正解趋于零。
\end{proposition}
\begin{proof}
定额模型右端是开口向下的抛物线，最大值为 $rK/4-h$。
当有两个根时，它在两根之间为正，在两根外为负。
所以 $x_-$ 两侧的运动都离开它，而 $x_+$ 两侧的运动都指向它。
唯一性禁止解穿过平衡点；单调有界解的极限必须是右端的零点，故得所述稳定性。

临界时右端为 $-(r/K)(x-K/2)^2$，从右侧趋近平衡点而从左侧离开，因而它是半稳定的。
超临界时令 $c=h-rK/4>0$，有 $\dot x\leq-c$，故不迟于时间 $x(0)/c$ 达到零。
达到零以后原方程会给出负数；种群模型在此已失去物理适用性。

比例模型右端为 $x[(r-q)-rx/K]$，按其零点两侧的符号作同样分析。
当 $q\geq r$ 时正解递减且非负，若极限为正，右端将有严格负的极限，
与收敛矛盾，所以极限为零。无捕捞情形的解由分离变量得到
\[
 x(t)=\frac{K}{1+(K/x_0-1)e^{-rt}},\qquad x_0>0.
\]
它也直接验证了 $K$ 对所有正初值的吸引性。
\end{proof}
